\documentclass[12pt]{article}
\makeatletter
\def\input@path{{../../}}
\makeatother
\usepackage{sbc-template}
\makeatletter
\renewcommand{\inst}[1]{}
\makeatother
\usepackage{graphicx,url}
\usepackage[utf8]{inputenc}
\usepackage[T1]{fontenc}
\usepackage[brazil]{babel}
\usepackage{longtable,booktabs,array}
\providecommand{\tightlist}{%%
  \setlength{\itemsep}{0pt}\setlength{\parskip}{0pt}}
\sloppy
\title{Previsao de demanda e otimizacao de estoque: por que este problema importa}
\author{}
\address{}
\begin{document}

\maketitle

\begin{resumo}
Este primeiro artigo abre a serie "QRC 101" com uma pergunta simples: por que previsao de demanda e um problema bom o bastante para sustentar uma obra inteira, da modelagem classica ate Quantum Reservoir Computing? A resposta tem duas camadas. A primeira e de negocio: erro de previsao vira ruptura, excesso de estoque, capital parado, compra emergencial e perda de margem. A segunda e metodologica: series de demanda sao temporais, sazonais, afetadas por promocoes, feriados e mudancas de comportamento, o que as torna um terreno natural para estudar modelos com memoria. Para concretizar a discussao, adotamos o dataset Favorita como fio condutor da serie e fixamos um recorte experimental reproduzivel: previsao diaria para \texttt{store\_nbr\ =\ 1} e \texttt{family\ =\ BEVERAGES}, com os ultimos 90 dias reservados para teste. Ao final do artigo, o leitor entende o problema de negocio, conhece os arquivos centrais do Favorita, sabe quais metricas guiam a avaliacao e enxerga como os sete artigos se conectam.
\end{resumo}

\section{1. Introducao}\label{1-introducao}

Poucos problemas de negocio sao tao simples de explicar e tao dificeis de resolver quanto previsao de demanda. A pergunta operacional parece modesta: quanto preciso comprar hoje para atender a demanda de amanha ou da proxima semana? Na pratica, essa pergunta conecta planejamento comercial, operacao logistica, capital de giro e experiencia do cliente. Se a empresa compra menos do que deveria, aparece ruptura e perda de venda. Se compra mais, aparece excesso de estoque, custo de armazenagem e risco de deterioracao ou obsolescencia.

Esse equilibrio e particularmente delicado no varejo alimentar, em que a demanda varia com calendario, promocao, comportamento regional e sazonalidade de curto e medio prazo. O problema nao e apenas prever um numero. O problema e prever uma serie temporal em um ambiente onde o passado importa, mas nao de maneira linear e nem sempre da mesma forma.

Essa combinacao faz do tema uma excelente porta de entrada para Reservoir Computing (RC) e, mais adiante, para Quantum Reservoir Computing (QRC). Em vez de comecar a serie pelo formalismo do reservatorio, comecamos pelo problema que justifica a modelagem temporal.

\section{2. Por que previsao de demanda importa para o negocio}\label{2-por-que-previsao-de-demanda-importa-para-o-negocio}

Em um contexto de estoque, o erro de previsao nao e abstrato. Ele gera consequencias concretas:

\begin{itemize}
\tightlist
\item
  subprevisao aumenta a chance de ruptura, perda de venda e piora do nivel de servico;
\item
  superprevisao aumenta estoque medio, capital imobilizado e desperdicio operacional;
\item
  previsoes instaveis pioram reabastecimento, negociacao com fornecedores e uso da malha logistica;
\item
  previsoes opacas dificultam explicar por que determinada decisao foi tomada.
\end{itemize}

Do ponto de vista didatico, esse problema e muito forte porque permite ligar tres niveis de analise:

\begin{enumerate}
\def\labelenumi{\arabic{enumi}.}
\tightlist
\item
  o nivel de negocio, em que a pergunta e "quanto venderemos?";
\item
  o nivel estatistico, em que a pergunta e "qual erro o modelo comete?";
\item
  o nivel operacional, em que a pergunta e "o erro ajuda ou atrapalha a decisao de estoque?".
\end{enumerate}

Ao longo da serie, essas tres perguntas serao mantidas juntas. Isso evita um erro comum em textos introdutorios: ensinar um modelo temporal sem explicar por que alguem se importaria com ele fora do laboratorio.

\section{3. Por que previsao de demanda e dificil}\label{3-por-que-previsao-de-demanda-e-dificil}

Series de demanda no varejo sao dificeis porque combinam padroes recorrentes com perturbacoes localizadas. Em um unico SKU ou familia de produtos, o valor observado em um dia pode refletir:

\begin{itemize}
\tightlist
\item
  efeito do dia da semana;
\item
  efeito do mes e de sazonalidade mais longa;
\item
  promocao local;
\item
  feriado nacional ou regional;
\item
  alteracao de fluxo na loja;
\item
  mudanca de comportamento do consumidor;
\item
  ruido operacional e eventos nao observados.
\end{itemize}

Isso significa que o problema nao pode ser reduzido a uma linha reta no tempo. Mesmo quando ha sazonalidade clara, ela pode ser modulada por contexto. Uma promocao de bebidas antes de um feriado nao tem o mesmo efeito de uma promocao identica em uma semana comum. Esse e justamente o tipo de dinamica em que memoria e nao linearidade se tornam relevantes.

Para a serie, isso cria um caminho pedagogico natural:

\begin{itemize}
\tightlist
\item
  primeiro, mostramos o problema com baselines simples;
\item
  depois, introduzimos RC como um mecanismo de representacao temporal;
\item
  por fim, mostramos como QRC preserva a mesma logica geral, mas troca o reservatorio classico por um reservatorio quantico.
\end{itemize}

\section{4. O caso Favorita}\label{4-o-caso-favorita}

O dataset adotado nesta obra e o da competicao "Corporacion Favorita Grocery Sales Forecasting". Ele e particularmente adequado para fins didaticos porque combina riqueza operacional com uma estrutura de tabelas facil de navegar.

Os arquivos centrais utilizados no projeto sao os seguintes:

\begin{longtable}[]{@{}lll@{}}
\toprule\noalign{}
Arquivo & Papel no negocio & Papel no projeto \\
\midrule\noalign{}
\endhead
\bottomrule\noalign{}
\endlastfoot
\texttt{train.csv} & historico diario de vendas & alvo principal da modelagem \\
\texttt{test.csv} & horizonte de previsao da competicao & referencia estrutural do problema \\
\texttt{stores.csv} & metadados das lojas & contexto institucional e geografico \\
\texttt{transactions.csv} & volume de transacoes por loja & proxy de fluxo e atividade \\
\texttt{holidays\_events.csv} & feriados e eventos & variaveis exogenas de calendario \\
\texttt{oil.csv} & serie de preco do petroleo & contexto macroeconomico da competicao \\
\texttt{sample\_submission.csv} & formato de submissao & referencia de saida do desafio \\
\end{longtable}

Mesmo sem usar toda a complexidade do dataset nos primeiros artigos, a estrutura do Favorita e valiosa porque permite escalar gradualmente o problema. Podemos comecar com uma serie unica e depois abrir caminho para configuracoes mais ricas, sempre sem trocar de dominio.

\section{5. Do problema de negocio ao problema de modelagem}\label{5-do-problema-de-negocio-ao-problema-de-modelagem}

Para transformar previsao de demanda em um problema reproduzivel de modelagem, precisamos fixar algumas escolhas.

\subsection{5.1 Alvo}\label{51-alvo}

O alvo principal da serie e a variavel de demanda diaria \texttt{y}, derivada de \texttt{unit\_sales} em um recorte controlado.

\subsection{5.2 Unidade de analise}\label{52-unidade-de-analise}

Em vez de atacar o dataset inteiro de uma vez, a serie fixa um recorte comum:

\begin{itemize}
\tightlist
\item
  \texttt{store\_nbr\ =\ 1}
\item
  \texttt{family\ =\ BEVERAGES}
\item
  frequencia diaria
\item
  ultimos \texttt{90} dias reservados para teste
\end{itemize}

Essa escolha tem tres vantagens:

\begin{enumerate}
\def\labelenumi{\arabic{enumi}.}
\tightlist
\item
  e simples o suficiente para ser reproduzida em qualquer maquina;
\item
  preserva relevancia de negocio, pois bebidas costumam responder a calendario e promocao;
\item
  cria um protocolo comum para baselines, modelos classicos, RC e QRC.
\end{enumerate}

\subsection{5.3 Variaveis exogenas}\label{53-variaveis-exogenas}

O projeto usa, em diferentes niveis de complexidade, informacoes como:

\begin{itemize}
\tightlist
\item
  \texttt{onpromotion};
\item
  componentes ciclicos de calendario;
\item
  informacoes agregadas de contexto, quando apropriado.
\end{itemize}

No pipeline sequencial usado por RC e QRC, o vetor de entrada parte de:

\begin{itemize}
\tightlist
\item
  demanda previa normalizada;
\item
  promocao normalizada;
\item
  \texttt{dow\_sin} e \texttt{dow\_cos};
\item
  \texttt{month\_sin} e \texttt{month\_cos};
\item
  indicador de fim de semana.
\end{itemize}

\subsection{5.4 Protocolo temporal}\label{54-protocolo-temporal}

Todos os modelos comparados na serie respeitam a mesma regra basica:

\begin{itemize}
\tightlist
\item
  treino no passado;
\item
  teste no bloco final dos ultimos 90 dias;
\item
  nenhuma informacao do futuro pode entrar no treino ou na engenharia de atributos.
\end{itemize}

Isso parece elementar, mas e justamente onde muitos experimentos de series temporais falham.

\section{6. O que sera considerado um bom resultado}\label{6-o-que-sera-considerado-um-bom-resultado}

Nesta serie, "bom resultado" nao quer dizer apenas erro pequeno. Quer dizer erro pequeno em um protocolo justo e com interpretacao util.

As metricas principais sao:

\begin{itemize}
\tightlist
\item
  \texttt{MAE}: media do erro absoluto, facil de interpretar em unidades de venda;
\item
  \texttt{RMSE}: penaliza mais fortemente erros grandes;
\item
  \texttt{WAPE}: relaciona erro absoluto ao volume total observado;
\item
  \texttt{sMAPE}: mede erro percentual de forma simetrica.
\end{itemize}

Cada metrica responde a uma pergunta diferente:

\begin{itemize}
\tightlist
\item
  \texttt{MAE}: o erro medio em unidades faz sentido operacional?
\item
  \texttt{RMSE}: o modelo esta cometendo erros muito grandes em alguns dias?
\item
  \texttt{WAPE}: o erro e alto ou baixo em relacao ao volume da serie?
\item
  \texttt{sMAPE}: o desempenho percentual e estavel em diferentes escalas?
\end{itemize}

Mais adiante, essas metricas serao complementadas por uma leitura de negocio: previsao nao serve apenas para "ganhar benchmark", mas para apoiar decisao de reposicao.

\section{7. Por que este problema e um bom fio condutor para RC e QRC}\label{7-por-que-este-problema-e-um-bom-fio-condutor-para-rc-e-qrc}

O caso Favorita permite ensinar toda a serie com continuidade:

\begin{itemize}
\tightlist
\item
  o problema nasce em uma dor real de negocio;
\item
  o sinal e claramente temporal;
\item
  ha espaço para baselines simples e modelos mais sofisticados;
\item
  o mesmo recorte pode ser usado para comparar familias de modelos muito diferentes.
\end{itemize}

Isso e importante porque um curso ou serie de artigos sobre QRC pode falhar de dois modos:

\begin{enumerate}
\def\labelenumi{\arabic{enumi}.}
\tightlist
\item
  ficar abstrato demais e perder o leitor;
\item
  pular para o quantico cedo demais e transformar tudo em curiosidade tecnica.
\end{enumerate}

Ao manter o mesmo problema de negocio do primeiro ao setimo artigo, a transicao para RC e depois para QRC deixa de parecer arbitraria.

\section{8. Roadmap da serie}\label{8-roadmap-da-serie}

O caminho proposto nesta obra e o seguinte:

\begin{enumerate}
\def\labelenumi{\arabic{enumi}.}
\tightlist
\item
  entender o problema de negocio e o recorte experimental;
\item
  introduzir Reservoir Computing com o minimo de formalismo necessario;
\item
  construir um primeiro pipeline reproduzivel;
\item
  explicar memoria, nao linearidade e capacidade de representacao;
\item
  estabelecer um protocolo justo de comparacao com baselines e IA classica;
\item
  construir a ponte conceitual entre RC classico, reservatorios fisicos e QRC;
\item
  fechar com um estudo aplicado de QRC no mesmo problema de demanda.
\end{enumerate}

Essa ordem foi escolhida para que o leitor nunca precise aceitar um modelo "por fe". Cada novo conceito entra para resolver uma pergunta aberta no artigo anterior.

\section{9. Conclusao}\label{9-conclusao}

Este primeiro artigo fez um trabalho fundamental: antes de falar em reservatorios, falamos em demanda, estoque e decisao. Escolhemos o Favorita como caso central porque ele sustenta tanto a motivacao de negocio quanto a complexidade metodologica necessaria para ensinar modelagem temporal de forma honesta.

Com o recorte experimental definido e as metricas estabelecidas, a serie pode agora dar o proximo passo: explicar Reservoir Computing nao como um truque exotico, mas como um paradigma de modelagem temporal.

\section{Entregaveis associados no repositorio}\label{entregaveis-associados-no-repositorio}

\begin{itemize}
\tightlist
\item
  visao geral da serie: \texttt{papers/README.md}
\item
  dataset do projeto: \texttt{dataset/raw/}
\item
  benchmark consolidado: \texttt{code/RESULTS.md}
\end{itemize}

\section{Referencias}\label{referencias}

\begin{enumerate}
\def\labelenumi{\arabic{enumi}.}
\tightlist
\item
  Hyndman, R. J., Athanasopoulos, G. Forecasting: Principles and Practice.
\item
  Kaggle. Corporacion Favorita Grocery Sales Forecasting.
\end{enumerate}

\section*{Resultados Computacionais Associados}

Os resultados computacionais associados a este artigo estao organizados na subpasta \texttt{computational\_results\_20260402\_222902/}, localizada dentro desta pasta \LaTeX. Esse diretorio contem o arquivo \texttt{REPORT.md}, tabelas em CSV e imagens comparativas apropriadas ao escopo do artigo.

\end{document}
